\documentclass{article}
\usepackage{graphicx}
\usepackage{amsmath}
\usepackage{amsfonts}
\usepackage{bookmark}
\usepackage{booktabs}
\usepackage{exercises}
\graphicspath{ {./images/10.4} }
\usepackage{indentfirst}
\usepackage{amsthm}
\usepackage{yhmath}
\usepackage{geometry}
\usepackage{placeins}
\usepackage{csquotes}
\usepackage{pgfplots}
\newcommand*{\oldneg}{\mathord{\sim}}
\newcommand\showdiv[1]{\overline{\smash{)}#1}}
\newcommand\Mydiv[1]{%
\smash{\raise.3ex\hbox{$\big)$}}\mkern-8mu
        \overline{\enspace\strut#1}}
\setlength{\tabcolsep}{5pt}
\theoremstyle{definition}
\newtheorem{definition}{Definition}

\setlength{\parindent}{.5in}
\setlength{\parskip}{12pt}
\renewcommand{\arraystretch}{2}
\geometry{margin=1in}
\newcommand{\?}{\stackrel{?}{=}}
\pgfplotsset{width=7cm,compat=1.18}\usepgfplotslibrary{statistics}

\title{10.5 Answers}
\begin{document}
\maketitle

\section*{3}
We are given a maximum of $(\frac{\pi}{4},3)$ and a minimum of $(\frac{3\pi}{4},-3)$. Thus the amplitude is $\frac{3+3}{2}=3$ and the vertical shift is $\frac{3-3}{2}=0$. Also, the period is
$2(\frac{3\pi}{4}-\frac{\pi}{4})=\pi$ and thus \[ \pi=\frac{2\pi}{b} \] \[ b=2 \] Because the $x$-coordinate of the maximum is a third of the minimum, we can model this as a sine function with no 
horizontal shift:
\[ y=3\sin 2x \]

\section*{4}
We are given a maximum of $(0,5)$ and a minimum of $(\frac{\pi}{4},-5)$. Thus the amplitude is $\frac{5+5}{2}=5$ and the vertical shift is $\frac{5-5}{2}=0$. Also, the period is 
$2(\frac{\pi}{4}-0)=\frac{\pi}{2}$ and thus \[ \frac{\pi}{2}=\frac{2\pi}{b} \] \[ b=4 \] Because the maximum is on the $y$-axis, we can model this as a cosine function with no horizontal shift or reflection:
\[ y=5\cos 4x \]

\section*{5}
We are given a maximum of $(2,6)$ and a minimum of $(0,2)$. Thus the amplitude is $\frac{6-2}{2}=2$ and the vertical shift is $\frac{6+2}{2}=4$. Also, the period is
$2(2-0)=4$ and thus \[ 4=\frac{2\pi}{b} \] \[ b=\frac{\pi}{2} \] We can therefore model this by a reflection of a cosine function with no horizontal shift:
\[ y=-2\cos\frac{\pi}{2}+4 \]

\section*{6}
We are given a maximum of $(\frac{3}{2},-1)$ and a minimum of $(\frac{1}{2},-3)$. Thus the amplitude is $\frac{-1+3}{2}=1$ and the vertical shift is $\frac{-1-3}{2}=-2$. Also, the period is
$2(\frac{3}{2}-\frac{1}{2})=2$ and thus \[ 2=\frac{2\pi}{b} \] \[ b=\pi \] Because the $x$-coordinate of the minimum is a third of the maximum, we can model this as a reflection of a sine function with 
no horizontal shift:
\[ y=-\sin\pi x-2 \]

\section*{9}
We are given that the maximum is $(\pi,6)$ and the minimum is $(3\pi,-6)$. Thus the amplitude is $\frac{6+6}{2}=6$ and the vertical shift is $\frac{6-6}{2}=0$. Also, we know that the period is
$2(3\pi-\pi)=4\pi$ and thus \[ 4\pi=\frac{2\pi}{b} \] \[ b=\frac{1}{2} \] We know that the $x$-coordinate of the maximum is a third of the minimum. Thus we one way we can model this is by a sine function 
with no reflection and no horizontal shift: \[ y=6\sin\frac{1}{2}x \]

\section*{10}
We are given that the maximum is $(0,4)$ and the minimum is $(\pi,-4)$. Thus the amplitude is $\frac{4+4}{2}=4$ and the vertical shift is $\frac{4-4}{2}=0$. Also, we know that the period is
$2(\pi-0)=2\pi$ and thus \[ 2\pi=\frac{2\pi}{b} \] \[ b=1 \] We know that the maximum occurs on the $y$-axis, meaning we can model this with a cosine functio with no horizontal shift:
\[ y=4\cos x\] 

\section*{11}
We are given that the maximum is $(\frac{\pi}{3},5)$ and the minimum is $(0,3)$. Thus the amplitude is $\frac{5-3}{2}=1$ and the vertical shift is $\frac{5+3}{2}=4$. Also, we know that the period is
$2(\frac{\pi}{3}-0)=\frac{2\pi}{3}$ and thus \[ \frac{2\pi}{3}=\frac{2\pi}{b} \] \[ b=3 \] Because the minimum occurs on the $y$-axis, we can model this by a reflection of a cosine function with no 
horizontal shift: \[ y=-\cos 3x+4 \]

\section*{12}
We are given that the maximum is $(\frac{\pi}{6},8)$ and the minimum is $(0,-6)$. Thus the amplitude is $\frac{8+6}{2}=7$ and the vertical shift is $\frac{8-6}{2}=1$. Also, we know that the period is
$2(\frac{\pi}{6}-0)=\frac{\pi}{3}$ and thus \[ \frac{\pi}{3}=\frac{2\pi}{b} \] \[ b=6 \] Because the minimum occurs on the $y$-axis, we can model this by a reflection of a cosine function with no horizontal shift:
\[ y=-7\cos 6x+1 \]

\section*{13}
We are given that the maximum is $(\frac{3\pi}{4},9)$ and the minimum is $(2\pi,5)$. Thus the amplitude is $\frac{9-5}{2}=2$ and the vertical shift is $\frac{9+5}{2}=7$. Also, we know that the period is
$2(2\pi-\frac{3\pi}{4})=2\frac{5\pi}{4}=\frac{5\pi}{2}$ and thus \[ \frac{5\pi}{2}=\frac{2\pi}{b} \] \[ b=\frac{4}{5} \] We can therefore model this is as a cosine function with no reflection and with a horizontal shift 
of $\frac{3\pi}{4}$: \[ y=2\cos\frac{4}{5}(x-\frac{3\pi}{4})+7 \]

\section*{14}
We are given that the maximum is $(0,5)$ and the minimum is $(6,-11)$. Thus the amplitude is $\frac{5+11}{2}=8$ and the vertical shift is $\frac{5-11}{2}=-3$. Also, we know that the period is
$2(6-0)=12$ and thus \[ 12=\frac{2\pi}{b} \] \[ b=\frac{\pi}{6} \] Because the maximum is at the $y$-axis, we can model this as a cosine function with no relfection and with a horiztonal shift of 0:
\[ y=8\cos\frac{\pi}{6}x-3 \]

\section*{15}
We are given that the maximum is $(0,0)$ and the minimum is $(4\pi,-4)$. Thus the amplitude is $\frac{0+4}{2}=2$ and the vertical shift is $\frac{0-4}{2}=-2$. Also, we know that the period is
$2(4\pi-0)=8\pi$ and thus \[ 8\pi=\frac{2\pi}{b} \] \[ b=\frac{1}{4} \] Because the maximum is at the $y$-axis, we can model this as a cosine function with no reflection and with a horizontal shift of 0:
\[ y=2\cos\frac{1}{4}x-2 \]

\section*{16}
We are given that the maximum is $(\frac{\pi}{3},-3)$ and the minimum is $(\frac{\pi}{12},-7)$. Thus the amplitude is $\frac{-3+7}{2}=2$ and the vertical shift is $\frac{-3-7}{2}=-5$. Also, we know that the period is 
$2(\frac{\pi}{3}-\frac{\pi}{12})=2(\frac{\pi}{4})=\frac{\pi}{2}$ and thus \[ \frac{\pi}{2}=\frac{2\pi}{b} \] \[ b=4 \] We can therefore model this as a cosine function with no reflection and with a horizontal shift of 
$\frac{\pi}{3}$: \[ y=2\cos 4(x-\frac{\pi}{3})-5 \]

\section*{17}
We are given that the maximum is $(\frac{2\pi}{3},0)$ and the minimum is $(0,-12)$. Thus the amplitude is $\frac{0+12}{2}=6$ and the vertical shift is $\frac{0-12}{2}=-6$. Also, we know that the period is
$2(\frac{2\pi}{3}-0)=\frac{4\pi}{3}$ and thus \[ \frac{4\pi}{3}=\frac{2\pi}{b} \] \[ b=\frac{3}{2} \] Because the minimum occurs on the $y$-axis, we can model this as a reflection of a cosine function with no 
horizontal shift: \[ y=-6\cos\frac{3}{2}x-6 \]

\section*{18}
The amplitude is $\frac{60-38}{2}=11$. Thus the answer is B. 

\section*{19}
We are given a maximum of $(6,15)$ and a minimum of $(0,5)$. Thus the amplitude is $\frac{15-5}{2}=5$ and the vertical shift is $\frac{15+5}{2}=10$. Also, we know that the period is
$2(6-0)=12$ and thus \[ 12=\frac{2\pi}{b} \] \[ b=\frac{\pi}{6} \] Because the minimum occurs on the $y$-axis, we can model this as a reflection of a cosine function with no reflection or horizontal shift:
\[ y=-5\cos\frac{\pi}{6}x+10 \]

\section*{20}
Suppose we are given a sinusoid whose $y$-intercept occurs at the minimum value of the function. Now, we know that cosine functions with no horiztonal shift are such that their maximum occurs at the $y$-intercept. 
Thus reflections of cosine functions are such that their minimum occurs at the $y$-intercept. Hence it would be convenient to model this as a cosine function.

\section*{21}
In Example 1, we modeled the graph with the equation $y=3\sin 4x+2$. Now, we know that
\[ \sin 4x=\cos(\frac{\pi}{2}-4x) \]
\[ =\cos(-(4x-\frac{\pi}{2})) \]
\[ =\cos(4x-\frac{\pi}{2}) \]
\[ =\cos 4(x-\frac{\pi}{8}) \]
meaning a cosine function to model the given equation is
\[ y=3\cos 4(x-\frac{\pi}{8})+2 \]

\section*{22}
We are given a sinusoid with a minimum at $(\frac{\pi}{2},3)$ and a maximum of $(\frac{\pi}{4},8)$. Thus the amplitude is $\frac{8-3}{2}=2.5$ and the vertical shift is 
$\frac{8+3}{2}=5.5$. Also, the period is $2(\frac{\pi}{2}-\frac{\pi}{4})=\frac{\pi}{2}$ and thus \[ \frac{\pi}{2}=\frac{2\pi}{b} \] \[ b=4 \] We can thus model this as a cosine function with a horiztonal shift of 
$\frac{\pi}{4}$:
\[ y=2.5\cos 4(x-\frac{\pi}{4})+5.5 \]
Now, we know that
\[ \cos 4(x-\frac{\pi}{4})=\sin(\frac{\pi}{2}-(4(x-\frac{\pi}{4}))) \]
\[ =\sin(\frac{\pi}{2}-4x+\pi) \]
\[ =\sin(\frac{3\pi}{2}-4x) \]
\[ =\sin(-(4x-\frac{3\pi}{2})) \]
\[ =-\sin(4x-\frac{3\pi}{2}) \]
\[ =-\sin 4(x-\frac{3\pi}{8}) \]
meaning an equivalent sine function for this is
\[ y=-2.5\sin 4(x-\frac{3\pi}{8})+5.5 \]

\section*{23}
We are given a maximum of $(\frac{1}{8},100)$ and a minimum of $(\frac{3}{8},-100)$, as well as a period of $0.5$, a vertical shift of $\frac{100-100}{2}=0$, and an amplitude of
$\frac{100+100}{2}=100$. Thus
\[ \frac{1}{2}=\frac{2\pi}{b} \] \[ b=4\pi \] Because the $x$-coordinate of the maximum is one third of the $x$-coordinate of the minimum, we can model this as a sine function with no horiztonal shift or reflection:
\[ y=100\sin 4\pi x \]

\section*{24}
We are given a minimum of $(0,52)$ and a maximum of $(6,84)$. Thus the amplitude is $\frac{84-52}{2}=16$ and the vertical shift is $\frac{84+52}{2}=68$. Also, the period is
$2(6-0)=12$ and thus \[ 12=\frac{2\pi}{b} \] \[ b=\frac{\pi}{6} \] Because the minimum occurs at the $y$-axis, we can model this as a reflection of a cosine function with no horizontal shift:
\[ y=-16\cos\frac{\pi}{6}x+68 \]

\section*{25}
We are given a radius of 5 feet, meaning the maximum is 25 feet above the ground, and that the minimum is 4 feet above the ground (when $t=0$). Thus the amplitude is 
$\frac{25-4}{2}=10.5$ and the vertical shift is $\frac{25+4}{2}=14.5$. We are also given that a complete turn is made every two seconds; hence the period is 2, meaning
\[ 2=\frac{2\pi}{b} \] \[ b=\pi \] Because the minimum is at the $y$-axis, we can model this as a reflection of a cosine function with no horizontal shift:
\[ y=-10.5\cos\pi x+14.5 \] 

\section*{26}
(a): We are given a minimum at $(0,3.5)$ and a maximum of $(6,16.5)$. Thus the amplitude is $\frac{16.5-3.5}{2}=6.5$ and the vertical shift is $\frac{3.5+16.5}{2}=10$. Also, the period is
$2(6-0)=12$ and thus \[ 12=\frac{2\pi}{b} \] \[ b=\frac{\pi}{6} \] Because the minimum occurs at the $y$-axis, this can be modeled by a reflection of a cosine function with no horizontal shift:
\[ f(x)=-6.5\cos\frac{\pi}{6}x+10 \] 


(b): We know that the period is 12. Thus the times at which high or low tide occurs is 0 hours, 6 hours, 12 hours, or 18 hours.


(c): The graph that shows the tide depths $d$ $t$ hours after 3 AM is just a composition of functions (let it be called $h(x)$). In particular, let $g(x)=t-3$. Then 
$h(x)=f(g(x))$, that is, $h(x)=f(t-3)$.

\end{document}
